% concatenated from seto_20260107.tex
\documentclass[12pt]{article}
\usepackage{amscd, enumerate, amsmath, amssymb}
\usepackage{tikz-cd}

\usepackage{amsthm}
\usepackage{url}
\renewcommand{\baselinestretch}{1.2}
\setlength{\oddsidemargin}{20pt}
\setlength{\topmargin}{-15pt}
%\setlength{\textwidth}{450pt}
%\setlength{\textheight}{660pt}
\setlength{\textwidth}{420pt}
\setlength{\textheight}{620pt}


\newtheorem{thm}{Theorem}[section]
\newtheorem{defn}[thm]{Definition}
\newtheorem{pf}{Proof}
\newtheorem{lem}[thm]{Lemma}
\newtheorem{rem}[thm]{Remark}
\newtheorem{ex}[thm]{Example}
\newtheorem{cor}[thm]{Corollary}
\newtheorem{prop}[thm]{Proposition}

\renewcommand{\thepf}{}
\renewcommand{\theenumi}{\roman{enumi}}
\renewcommand{\labelenumi}{(\theenumi)}
\renewcommand{\refname}{References}

\newtheorem{acknowledgment}{Acknowledgment}
\renewcommand{\theacknowledgment}{}


\numberwithin{equation}{section}
%\setcounter{section}{1}



\newcommand{\N}{{\mathcal N}}
\newcommand{\cD}{{\mathcal D}}
\newcommand{\cB}{{\mathcal B}}
\newcommand{\cR}{{\mathcal R}}
\newcommand{\E}{{\mathcal E}}
\newcommand{\F}{{\mathcal F}}
\newcommand{\M}{{\mathcal M}}
\newcommand{\cL}{{\mathcal L}}
\newcommand{\cP}{{\mathcal P}}

\newcommand{\C}{{\mathbb C}}
\newcommand{\R}{{\mathbb R}}
\newcommand{\D}{{\mathbb D}}
\newcommand{\bT}{{\mathbb T}}
\newcommand{\Z}{{\mathbb Z}}
\newcommand{\bH}{{\mathbb H}}

\newcommand{\cS}{{\mathcal S}}

\newcommand{\bI}{{\mathbb I}}
\newcommand{\bJ}{{\mathbb J}}

\newcommand{\cH}{{\mathcal H}}
\newcommand{\cK}{{\mathcal K}}
\newcommand{\cG}{{\mathcal G}}
\newcommand{\ran}{\operatorname{ran}}
\newcommand{\Span}{\operatorname{span}}
\newcommand{\tr}{\operatorname{tr}}
\newcommand{\Real}{\operatorname{Re}}

\newcommand{\la}{\langle}
\newcommand{\ra}{\rangle}
\newcommand{\p}{\partial}
\newcommand{\lam}{\lambda}




\begin{document}
\title{L\"{o}wner equations and de Branges--Rovnyak spaces}
\author{Michio Seto\\
Department of Mathematics, \\
National Defense Academy of Japan\\
{\small 
{\it E-mail address}: {\tt mseto@nda.ac.jp}}\\
} 
\maketitle

\begin{abstract} 
We study de Branges--Rovnyak spaces parametrized by L\"{o}wner equations 
after de Branges and Ghosechowdhury.  
The purpose of this paper is to propose a new approach based on the L\"{o}wner theory 
to the structure theory of de Branges--Rovnyak spaces. 
As the main results, three solutions to the problem ``Find concrete elements in de Branges--Rovnyak spaces'' discussed by 
Sarason, Fricain--Mashreghi, Fricain--Hartmann--Ross and {\L}anucha--Nowak are given. 
In particular, we also deal with applications of the chordal  L\"{o}wner theory to the de Branges--Rovnyak spaces. 
\end{abstract}

\begin{center}
2020 Mathematical Subject Classification: Primary 30C40; Secondary 46E22\\
keywords: L\"{o}wner equation, de Branges--Rovnyak space, reproducing kernel Hilbert space
\end{center}

\section{Introduction}
In de Branges' original proof of the Bieberbach conjecture, 
the fantastic blend of the L\"{o}wner theory and the de Branges--Rovnyak theory is one of the highlights (see de Branges~\cite{de Branges2} and Vasyunin--Nikol'ski\u{\i}~\cite{VN1, VN2}). 
Although this functional analysis aspect of de Branges' proof cannot be seen in 
the widely circulated proof (for example, Conway~\cite{Conway} and Rosenblum--Rovnyak~\cite{RR}),  
Ghosechowdhury continued to search for 
interaction between 
the L\"{o}wner theory and the de Branges--Rovnyak theory in \cite{G1, G2}, 
and obtained a Fourier type expansion theorem induced 
by the L\"{o}wner equation. 
More precisely, he showed continuous resolution 
of de Branges--Rovnyak spaces induced by univalent (i.e. holomorphic and injective) self-maps of the open unit disk.  
Now, we would like to point out that 
Ghosechowdhury's theorem leads us to a new approach to the problem 
\begin{center}
``\textit{Find concrete 
elements in de Branges--Rovnyak spaces}'' 
\end{center}
discussed by Sarason~\cite{Sarason}, 
Fricain--Mashreghi~\cite{FM}, Fricain--Hartmann--Ross~\cite{FHR} and {\L}anucha--
Nowak~\cite{LN}, which we will call the inhabitant problem. 
In fact, 
we show that
\[
\log \frac{1-B_t(z)}{1-z}
\]
belongs to the de Branges--Rovnyak space induced by $B_t$, 
where $B_t(z)=B(t,a,z)$ $(0\leq a \leq t)$ is the L\"{o}wner semigroup 
associated with the Koebe functions
\[
f(t,z)=\frac{e^tz}{(1-z)^2}\quad (t\geq 0)
\]
(Corollary \ref{cor:3-7}). 
Thus, our method led by Ghosechowdhury's theorem gives an answer to Question 6.6 in \cite{FHR}. 
However, it is probably because Ghosechowdhury's paper was written with the language of the de Branges school, 
it is hard to say that his work is widely known.  
The author's early motivation for this project was to reconstruct his work 
in the style of Ando~\cite{Ando}. 
In order to rearrange his proof, we adopt direct integrals and integral operators. 
Further, in our framework, 
applications of the chordal L\"{o}wner theory (see Schlei{\ss}inger~\cite{Sch}) to de Branges--Rovnyak spaces are discussed, too. 
The author hopes that this paper  
would make Ghosechowdhury's work accessible 
to Hilbert space operator theorists. 

This paper is organized as follows. 
Sections 2 and 3 are the preliminary parts, where direct integrals of reproducing kernel Hilbert spaces and 
integral operators defined on those spaces are introduced. 
In Section 4, 
we prove Ghosechowdhury's theorem in our style, and give a solution to the inhabitant problem.  
In Section 5, toward Section 6, we introduce auxiliary spaces defined by Pick functions. 
In Section 6, 
we discuss applications of the chordal L\"{o}wner theory to the de Branges--Rovnyak theory. 
In Corollaries \ref{cor:5-5} and \ref{cor:6-4}, further solutions to the inhabitant problem are obtained. 

\section{Direct integrals}\label{sec:1}

In this section, we arrange the theory of direct integrals to fit the L\"{o}wner theory. 
\subsection{Setting}
Let $\Omega$ be a domain in the complex plane $\C$, and let $(X,\sigma)$ be a measure space.  
We consider a family $\{k(x,z,w)\}_{x\in X}$ consisting of kernel functions (positive semidefinite functions) on $\Omega\times \Omega$. 
Let $C$ be any compact subset of $\Omega$. 
Then, we will write 
\[
M_C(x)=\sup_{z\in C}k(x,z,z).
\] 
Throughout this section, we assume the following four conditions: 
\begin{enumerate}
\item[(A)] $k(x,z,w)$ is measurable with respect to $x$ when $z$ and $w$ are fixed, 
\item[(B)] $k(x,z,w)$ is holomorphic with respect to $z$ when $x$ and $w$ are fixed, 
\item[(C)] $M_C(x)$ is integrable with respect to $\sigma$ for any $C$. In particular, $k(x,z,z)$ is integrable for any fixed $z$. 
\item[(D)] $M_C(x)$ is finite for any $C$ when $x$ is fixed. 
\end{enumerate}
In what follows, 
to avoid confusion, complex variables will be denoted such as $z$ and $w$, 
and specific points in $\Omega$ will be denoted such as $\lam$ and $\mu$. 
We will write $k_{\lam}(x,z)=k(x,z,\lam)$ for any $\lam$ in $\Omega$. 

\subsection{$\cH(x)$}
Let $\cH(x)$ denote the reproducing kernel Hilbert space generated 
by the kernel function $k(x,z,w)$ for a point $x$ in $X$. 

\begin{prop}\label{prop:4-1} 
Let $x$ be a point $x$ in $X$. 
Then, $\cH(x)$ is a reproducing kernel Hilbert space consisting of holomorphic functions on $\Omega$. 
\end{prop}

\begin{proof}
For any function $f(x,z)$ in $\cH(x)$, 
there exists a sequence $\{f_n(x,z)\}_n$ consisting of linear combinations of reproducing kernels of $\cH(x)$ 
such that 
\[
\lim_{n\to \infty}\|f_n(x,\cdot)-f(x,\cdot)\|_{\cH(x)}=0.
\]
Then, for any $\lam$ in $\Omega$, we have that 
\begin{align*}
|f_n(x,\lam)-f(x,\lam)|
&=|\la f_n(x,\cdot)-f(x,\cdot),k_{\lam}(x,\cdot) \ra_{\cH(x)}|\\
&\leq \|f_n(x,\cdot)-f(x,\cdot)\|_{\cH(x)}\|k_{\lam}(x,\cdot)\|_{\cH(x)}\\
&\to 0\quad(n\to \infty).
\end{align*}
It follows from Assumption (D) that  
this convergence is uniform on any compact subset of $\Omega$.   
Hence, by Assumption (B), $f(x,z)$ is holomorphic in $\Omega$ with respect to variable $z$. 
\end{proof}


\subsection{Riesz--Fischer type Lemma}
Let $L^{\infty}(X)$ denote the set of all bounded measurable functions on $X$. 
We need the following auxiliary space consisting of functions with two variables $x$ and $z$: 
\[
\cK=\left\{f:f(x,z)=\sum_{j=1}^nf_j(x)k_{\lam_j}(x,z),\ f_j\in L^{\infty}(X),\ \lam_j\in \Omega,\ n\geq 1\right\}.
\]
Then, $\cK$ is a vector space and invariant under multiplication by functions in $L^{\infty}(X)$. 
We shall mention the following rather trivial fact as a proposition. 
\begin{prop}
Let $f$ be any function in $\cK$. 
Then, $\|f(x,\cdot)\|_{\cH(x)}$ is a square integrable function with variable $x$. 
\end{prop}

\begin{proof}
We set $f(x,z)=\sum_{j=1}^nf_j(x)k_{\lam_j}(x,z)$. 
Then, we have that 
\begin{align*}
\|f(x,\cdot)\|_{\cH(x)}^2&=\sum_{i,j=1}^n\overline{f_i(x)}f_j(x)\la k_{\lam_j}(x,\cdot),k_{\lam_i}(x,\cdot) \ra_{\cH(x)}\\
&=\sum_{i,j=1}^n\overline{f_i(x)}f_j(x)k(x,\lam_i,\lam_j). 
\end{align*}
Hence, by Assumption (A), $\|f(x,\cdot)\|_{\cH(x)}^2$ is measurable.  
Moreover, by Assumption (C), $\|k_{\lam_j}(x,\cdot)\|_{\cH(x)}$ is square integrable. 
Hence, we have that 
\begin{align*}
\int_X \|f(x,\cdot)\|_{\cH(x)}^2\ d\sigma(x)
&\leq \int_X\sum_{i,j=1}^n\left|\overline{f_i(x)}f_j(x)k(x,\lam_i,\lam_j)\right|\ d\sigma(x)\\
&\leq \sum_{i,j=1}^n\|f_i\|_{\infty}\|f_j\|_{\infty}\int_X\|k_{\lam_i}\|_{\cH(x)}\|k_{\lam_j}\|_{\cH(x)}\ d\sigma(x)\\
&<\infty.
\end{align*}
This concludes the proof.
\end{proof}

We need the following Riesz--Fischer type lemma. 

\begin{lem}\label{lem:4-1}
Let $\{f_n\}_n$ be a sequence in $\cK$. 
If 
\[
\int_X\|f_n(x,\cdot)-f_m(x,\cdot)\|_{\cH(x)}^2\ d\sigma(x)\to 0\quad(n,m\to \infty),
\] 
then, 
for almost all $x$,   
there exists a function $f(x, z)$ in $\cH(x)$ such that 
\[
\int_X\|f(x,\cdot)\|_{\cH(x)}^2\ d\sigma(x)<\infty
\]
and
\[
\int_X\|f_n(x,\cdot)-f(x,\cdot)\|_{\cH(x)}^2\ d\sigma(x)\to 0\quad(n\to \infty).
\] 
\end{lem}

\begin{proof}
Since this lemma is essentially the same as the Riesz--Fischer theorem, 
we give a sketch of the proof. 
It follows from the assumption that 
there exists a subsequence $\{f_{n_k}\}_k$ of $\{f_n\}_n$ such that 
\[
\int_X \|f_{n_{k+1}}(x,\cdot)-f_{n_k}(x,\cdot)\|_{\cH(x)}^2\ d\sigma(x)<\frac{1}{2^{2k}}. 
\]
Then, since 
\[
g(x)=\|f_{n_1}(x,\cdot)\|_{\cH(x)}+\sum_{k=1}^{\infty}\|f_{n_{k+1}}(x,\cdot)-f_{n_k}(x,\cdot)\|_{\cH(x)} 
\]
is square integrable by the monotone convergence theorem. 
In particular, $g(x)$ is finite for almost all $x$ in $X$. 
Hence,
we have that
\[
\|f_{n_{\ell}}(x,\cdot)-f_{n_k}(x,\cdot)\|_{\cH(x)}
\leq \sum_{j=k}^{\ell-1}\|f_{n_{j+1}}(x,\cdot)-f_{n_j}(x,\cdot)\|_{\cH(x)}\to 0\quad(k,\ell\to\infty)
\]
for almost all $x$ in $X$. 
Therefore, 
for almost all $x$,   
there exists a function $f(x, z)$ in $\cH(x)$ 
such that 
\[
\lim_{k\to \infty}\|f_{n_k}(x,\cdot)-f(x,\cdot)\|_{\cH(x)}=0. 
\]
Then, 
since $\|f_{n_k}(x,\cdot)\|_{\cH(x)}\leq g(x)$, 
by Lebesgue's dominated convergence theorem, we have 
\[
\int_X\|f(x,\cdot)\|_{\cH(x)}^2\ d\sigma(x)<\infty.
\]
Further, 
by Fatou's lemma, we have 
\[
\int_X\|f_n(x,\cdot)-f(x,\cdot)\|_{\cH(x)}^2\ d\sigma(x)\to 0\quad(n\to \infty).
\] 
\end{proof}


\subsection{$L^2(\cH(x))$}
Let $L^2(\cH(x))$ denote the set of all functions obtained by the manner of Lemma \ref{lem:4-1}. 
Setting 
\[
\|f\|_{L^2(\cH(x))}=\left(
\int_X \|f(x,\cdot)\|_{\cH(x)}^2\ d\sigma(x)
\right)^{1/2}, 
\]
we may write  
\[
L^2(\cH(x))=\left\{f:\exists f_n\in \cK\ \mbox{s.t.}\ 
\lim_{n\to \infty}\|f_n-f\|_{L^2(\cH(x))}=0\right\}. 
\]
In what follows, two functions 
$f$ and $g$ in $L^2(\cH(x))$ will be identified if 
$\|f-g\|_{L^2(\cH(x))}=0$.
We summarize the basic properties on $L^2(\cH(x))$ as follows: 

\begin{thm}\label{thm:2-4}
$L^2(\cH(x))$ is an $L^{\infty}$-Hilbert module consisting of functions with two variables $x$ and $z$. 
Moreover, let $f$ be any function in $L^2(\cH(x))$. Then, 
$f(x,z)$ is a measurable function with variable $x$ when $z$ is fixed and 
$f(x,z)$ is a holomorphic function on $\Omega$ with variable $z$ for almost all $x$. 
\end{thm}

\begin{proof}
The proof of the first half is the same as that of Lemma \ref{lem:4-1}. 
We give the proof of the second half. 
Let $\{f_n\}_n$ be a sequence in $\cK$ such that 
\[
\lim_{n\to \infty}\|f_n-f\|_{L^2(\cH(x))}=0. 
\]   
Moreover, without loss of generality, we may assume that 
\[
\lim_{n\to \infty}\|f_n(x,\cdot)-f(x,\cdot)\|_{\cH(x)}=0
\]
for almost all $x$ in $X$.
Then, for any $\lam$ in $\Omega$, we have that
\[
|f_n(x,\lam)-f(x,\lam)|
\leq \|f_n(x,\cdot)-f(x,\cdot)\|_{\cH(x)}\|k_{\lam}(x,\cdot)\|_{\cH(x)}
\to 0\quad(n\to \infty).
\]
By Assumptions (A), (B) and (D), we have the conclusion. 
\end{proof}

\begin{ex}\label{ex:2-5}\rm 
Suppose that $h(x,z)$ is a nonvanishing function on $X\times \Omega$. 
We shall consider kernel function $k(x,z,w)=\overline{h(x,w)}h(x,z)$. Then, 
it is easy to see that 
\[
\cK=\{f:f(x,z)=g(x)h(x,z),\ g\in L^{\infty}(X)\}
\] 
and 
\[
\|f\|_{L^2(\cH(x))}^2=\int_X|g(x)|^2\|h(x,\cdot)\|_{\cH(x)}^2\ d\sigma(x)
\]
for any $f(x,z)=g(x)h(x,z)$ in $\cK$. 
Moreover, since
\[
|h(x,w)|^2=k(x,w,w)=\left\|\overline{h(x,w)}h(x,\cdot)\right\|_{\cH(x)}^2=|h(x,w)|^2\|h(x,\cdot)\|_{\cH(x)}^2, 
\]
we have that $\|h(x,\cdot)\|_{\cH(x)}^2=1$. 
Hence, if $\{f_n\}_n$ is a Cauchy sequence in $L^2(\cH(x))$, then 
so is $\{g_n\}_n$ in $L^2(X)$, where we set $f_n(x,z)=g_n(x)h(x,z)$.  
This observation concludes that 
\[
L^2(\cH(x))=\{f:f(x,z)=g(x)h(x,z),\ g \in L^2(X)\}.
\]
\end{ex}

\section{Integral operators}

Let $\{B(x,z)\}_{x\in X}$ be a family consisting of holomorphic functions on $\Omega$ satisfying
\begin{enumerate}
\item[(E)] $\|B\|_{\infty,C}=\sup_{x\in X,z \in C}|B(x,z)|$ is finite for any compact subset $C$ of $\Omega$. 
\end{enumerate}
In this section, we consider the integral operator defined as follows:
\[
\bI_B: f\mapsto \int_XB(x,z)f(x,z)\ d\sigma(x)\quad (f\in L^2(\cH(x))).
\]


\subsection{Basic properties}

First, we note that $(\bI_B f)(\lam)$ is finite for any $\lam$ in $\Omega$. 
Indeed, by Assumption (E)  we have that
\begin{align*}
\int_X|B(x,\lam)f(x,\lam)|\ d\sigma(x)
&\leq \|B\|_{\infty,\{\lam\}}\int_X|f(x,\lam)|\ d\sigma(x)\\
&\leq \|B\|_{\infty,\{\lam\}}\int_X\|f(x,\cdot)\|_{\cH(x)}\|k_{\lam}(x,\cdot)\|_{\cH(x)}\ d\sigma(x)\\
&\leq \|B\|_{\infty,\{\lam\}}\|f\|_{L^2(\cH(x))}\|k_{\lam}\|_{L^2(\cH(x))}.
\end{align*}
Moreover, we have that
\begin{equation}\label{eq:2-1}
|(\bI_B f)(\lam)-(\bI_B g)(\lam)|\leq \|B\|_{\infty,\{\lam\}}\|f-g\|_{L^2(\cH(x))}\|k_{\lam}\|_{L^2(\cH(x))}
\end{equation}
for any $f$ and $g$ in $L^2(\cH(x))$. 

\begin{lem}\label{lem:1-2-0}
For any $\lam$ in $\Omega$, $\bI_B k_{\lam}$ is a holomorphic function on $\Omega$. 
\end{lem}

\begin{proof}
Let $a$ be a point in $\Omega$ and 
let $c_j(x,\lam)$ denote the $j$-th coefficient of the power series expansion centered at $a$ of $B(x,z)k_{\lam}(x,z)$. 
Then, 
choosing appropriate $R>0$, 
we have that 
\begin{align*}
|c_j(x,\lam)|
&\leq \frac{1}{2\pi}\int_0^{2\pi}\frac{|B(x,a+Re^{i\theta})k_{\lam}(x,a+Re^{i\theta})|}{R^j}\ d\theta\\
&\leq \frac{\|B\|_{\infty,C}}{2\pi}\int_0^{2\pi}\frac{\|k_{a+Re^{i\theta}}(x,\cdot)\|_{\cH(x)}\|k_{\lam}(x,\cdot)\|_{\cH(x)}}{R^j}\ d\theta\\
&\leq \frac{\|B\|_{\infty,C}\sqrt{M_C(x)}\|k_{\lam}(x,\cdot)\|_{\cH(x)}}{R^j},
\end{align*}
where $C$ is chosen as $C=\{a+Re^{i\theta}:0\leq \theta <2\pi\}$.
Hence, if $|z-a|<r<R$, then
\begin{align*}
\int_X \left(\sum_{j=0}^{\infty}|c_j(x,\lam)(z-a)^j|\right)\ d\sigma(x)
&\leq \int_X \left(\sum_{j=0}^{\infty}\frac{\|B\|_{\infty,C}\sqrt{M_C(x)}\|k_{\lam}(x,\cdot)\|_{\cH(x)}}{R^j}r^j\right)\ d\sigma(x)\\
&=\frac{R}{R-r}\|B\|_{\infty,C}\left(\int_X M_C(x)\ d\sigma(x)\right)^{1/2}\|k_{\lam}\|_{L^2(\cH(x))}\\
&<\infty
\end{align*}
by Assumption (C). 
Hence, by Fubini's theorem, 
\begin{align*}
(\bI_B k_{\lam})(z)
&=\int_X B(x,z)k_{\lam}(x,z)\ d\sigma(x)\\
&=\int_X \left(\sum_{j=0}^{\infty}c_j(x,\lam)(z-a)^j\right)\ d\sigma(x)\\
&=\sum_{j=0}^{\infty}\left(\int_X c_j(x,\lam)\ d\sigma(x) \right)(z-a)^j
\end{align*}
is holomorphic at $a$. 
\end{proof}

\begin{lem}\label{lem:2-2}
The range space of $\bI_B$ is a subspace of the vector space consisting of holomorphic functions on $\Omega$ 
and $\ker \bI_B$ is closed. 
\end{lem}

\begin{proof}
Let $f$ be any function in $L^2(\cH(x))$. Then, by definition, there exists a sequence $\{f_n\}_n$ in $\cK$ such that 
\[
\lim_{n\to \infty}\|f_n-f\|_{L^2(\cH(x))}=0.
\]
Then, by (\ref{eq:2-1}), $(\bI_B f_n)(\lam)$ converges to $(\bI_B f)(\lam)$ for every $\lam$ in $\Omega$. 
Moreover, by Assumptions (D) and (E), this convergence is uniform on any compact subset of $\Omega$. 
Hence,  by Lemma \ref{lem:1-2-0}, $(\bI_B f)(z)$ is a holomorphic function on $\Omega$. 
Next, 
let $f$ be a function in the closure of $\ker \bI_B$ and let $\{f_n\}_n$ be a sequence in $\ker \bI_B$ such that 
\[
\lim_{n\to \infty}\|f_n-f\|_{L^2(\cH(x))}=0.
\]
Then, 
by (\ref{eq:2-1}), we have that
\begin{align*}
|(\bI_B f)(\lam)|
&=|(\bI_B f)(\lam)-(\bI_B f_n)(\lam)|\\
&\leq \|B\|_{\infty,\{\lam\}} \|f-f_n\|_{L^2(\cH(x))}\|k_{\lam}\|_{L^2(\cH(x))}\\
&\to 0\quad(n\to \infty)
\end{align*}
for any $\lam$ in $\Omega$. 
Hence, $\ker \bI$ is closed. 
\end{proof}

By Lemma \ref{lem:2-2} and the pull-back construction, 
we can endow the range space of $\bI_B$ with a Hilbert space structure as follows: 
\begin{align*}
\M(\bI_B; L^2(\cH(x)))
&=\M(\bI_B)=(\ran \bI_B,\ \|\cdot\|_{\M(\bI_B)}),\\
\|\bI_B f\|_{\M(\bI_B)}
&=\|P_{(\ker \bI_B)^{\perp}}f\|_{L^2(\cH(x))},
\end{align*}
where $(\ker \bI_B)^{\perp}$ denotes the orthogonal complement of $\ker \bI_B$ in $L^2(\cH(x))$. 
\begin{thm}\label{thm:2-2}
$\M(\bI_B)$ is a reproducing kernel Hilbert space and 
its reproducing kernel is 
\[
\left(\bI_B \overline{B(\cdot,\lam)}k_{\lam}\right)(z)=\int_X\overline{B(x,\lam)}B(x,z)k_{\lam}(x,z)\ d\sigma(x).
\]
In particular, for any function $F$ in $\M(\bI_B)$, there exists some function $f$ in $L^2(\cH(x))$ such that 
\[
F(z)=\int_XB(x,z)f(x,z)\ d\sigma(x).
\]
\end{thm}

\begin{proof}
We set $h_{\lam}(x,z)=\overline{B(x,\lam)}k_{\lam}(x,z)$. 
If $g$ belongs to $\ker \bI_B$, then
\begin{align*}
\la g, h_{\lam} \ra_{L^2(\cH(x))}
&=\int_{X}\la g(x,\cdot), h_{\lam}(x,\cdot) \ra_{\cH(x))}\ d\sigma(x)\\
&=\int_{X}B(x,\lam)g(x,\lam)\ d\sigma(x)\\
&=(\bI_B g)(\lam)\\
&=0. 
\end{align*}
This concludes that $h_{\lam}$ belongs to $(\ker \bI_B)^{\perp}$. 
Hence, for any function $f$ in $L^2(\cH(x))$, we have that 
\begin{align*}
\la \bI_B f, \bI_B h_{\lam} \ra_{\M(\bI)}
&=\la f,  h_{\lam}\ra_{L^2(\cH(x))}\\
&=\int_{X}\la f(x,\cdot),  h_{\lam}(x,\cdot)\ra_{\cH(x)}\ d\sigma(x)\\
&=\int_{X}B(x,\lam)f(x,\lam)\ d\sigma(x)\\
&=(\bI_B f)(\lam).
\end{align*}
Therefore, 
$\bI_B h_{\lam}$ is the reproducing kernel of $\M(\bI_B)$. 
\end{proof}


\begin{rem}\rm 

In de Branges~\cite{de Branges2}, the following type of integral operator was studied: 
\[
\bJ_B: f\mapsto \int_Xf(x,B(x,z))\ d\sigma(x)\quad (f\in L^2(\cH(x))). 
\]
In our framework, 
under suitable conditions on $k$ and $B$, 
it is shown that 
$(\bJ_Bf)(\lam)$ is finite for any $\lam$ in $\Omega$ and 
\[
|(\bJ_B f)(\lam)-(\bJ_B g)(\lam)|\leq \|f-g\|_{L^2(\cH(x))}
\left( \int_Xk(x,B(x,\lam),B(x,\lam))\ d\sigma(x)\right)^{1/2}
\]
for any $f$ and $g$ in $L^2(\cH(x))$. It follows from this estimate that 
$\M(\bJ_B)$ is a reproducing kernel Hilbert space with kernel function 
\[
\int_Xk(x,B(x,z),B(x,\lam))\ d\sigma(x). 
\]
\end{rem}


\subsection{Examples}


\begin{ex}[Paley--Wiener spaces]\rm 
Let $A$ be a positive constant. 
We consider kernel function 
\[
\frac{\sin (2\pi A(z-\overline{w}))}{\pi(z-\overline{w})}
\]
defined on $\C\times \C$. 
Then, the reproducing kernel Hilbert space generated by this kernel function  
is a Paley--Wiener space, which will be dented as $PW_A$. 
In Theorem \ref{thm:2-2}, 
we consider the case where $k(t,z,w)=1$, $X=[-A,A]$, $d\sigma=dt$ and $B(t,z)=e^{-2\pi i zt}$. Then, 
$\cH(t)=\C$, 
$L^2(\cH(t))=L^2([-A,A])$. 
Hence, 
\[
\bI_B: 
f(t,z) \mapsto F(z)=\int_{-A}^A B(t,z)f(t)\ dt \quad (f\in L^2([-A,A])).
\]
is the time-limited Fourier transform. 
Moreover, 
$\M(\bI_B; L^2([-A,A]))=PW_A$ as reproducing kernel Hilbert spaces. 
To see this, by Theorem \ref{thm:2-2}, it suffices to show that $\M(\bI_B)$ has the same reproducing kernel as that of $PW_A$.   
Setting $h_{\lam}(t,z)=\overline{B(t,\lam )}k_{\lam}(t,z)$, we have that 
\[
(\bI_B h_{\lam})(z)
=\int_{-A}^AB(t,z)h_{\lam}(t,z)\ dt\\
=\int_{-A}^Ae^{-2\pi i t(z-\overline{\lam})}\ dt\\
=\frac{\sin (2\pi A(z-\overline{\lam}))}{\pi(z-\overline{\lam})}. 
\]
Thus, we have $\M(\bI_B)=PW_A$. 
\end{ex}

\begin{ex}[Herglotz spaces]\label{sec:1-1-6}\rm 
Let $\varphi$ be a Herglotz function, that is, 
$\varphi$ is a holomorphic function on $\D$ and $\operatorname{Re}\varphi(z)\geq 0$ and $\varphi(0)=1$.  
Then, 
$\varphi$ has the following integral representation: 
\[
\varphi(z)=\int_{\bT}\frac{1+\xi z}{1-\xi z}\ d\mu(\xi),
\]
where $\mu$ is a probability measure on the unit circle $\bT$. 
The integrand 
\[
\varphi_{\xi}(z)=\frac{1+\xi z}{1-\xi z}
\]
is a typical example of a Herglotz function. 
Then, since
\begin{equation}\label{eq:3-3}
\frac{\overline{\varphi_{\xi}(w)}+\varphi_{\xi}(z)}{1-\overline{w}z}=
\frac{2}{(1-\overline{\xi w})(1-\xi z)},
\end{equation}
we have the following well-known fact.
\begin{prop}\label{prop:1-3}
Let $\varphi$ be a Herglotz function. Then, 
\[
\dfrac{\overline{\varphi(w)}+\varphi(z)}{1-\overline{w}z}
\]
is a kernel function on $\D\times \D$.
\end{prop}

The reproducing kernel Hilbert space generated by the kernel function in Proposition \ref{prop:1-3} 
is called a Herglotz space, which will be dented as $\cL(\varphi)$. 
Herglotz spaces were introduced and studied by de Branges--Rovnyak~\cite{dBR} 
(see also Ghosechowdhury~\cite{G1,G2} and Shulman~\cite{S}). 
Next, we consider the integral operator defined as follows: 
\[
\bI: 
f(\xi,z) \mapsto F(z)=\int_{\bT} f(\xi,z)\ d\mu(\xi) \quad (f\in L^2(\cL(\varphi_{\xi}))),
\]
which is obtained by setting $B\equiv 1$. 
Then,  
$\M(\bI; L^2(\cL(\varphi_{\xi})))=\cL(\varphi)$ as reproducing kernel Hilbert spaces. 
To see this, by Theorem \ref{thm:2-2}, 
it suffices to show that $\M(\bI)$ has the same reproducing kernel as that of $\cL(\varphi)$.   
Setting $k_{\lam}(\xi,z)=k(\xi,z,\lam)$, we have that 
\[
(\bI k_{\lam})(z)=\int_{\bT}k_{\lam}(\xi,z)\ d\mu(\xi)
=\int_{\bT} \frac{\overline{\varphi_{\xi}(\lam)}+\varphi_{\xi}(z)}{1-\overline{\lam} z}\ d\mu(\xi)
=\frac{\overline{\varphi(\lam)}+\varphi(z)}{1-\overline{\lam} z}. 
\]
Thus, we have $\M(\bI)=\cL(\varphi)$. 
\end{ex}


\section{L\"{o}wner expansions}
Let $\{B_{rs}(z)\}_{0\leq r\leq s}$ be a L\"{o}wner semigroup, that is, 
each $B_{rs}(z)$ is a univalent (i.e. holomorphic and injective) self-map of $\D$ satisfying 
the nonlinear L\"{o}wner equation
\begin{equation}\label{eq:5-1}
\frac{\p B_{rs}}{\p s}(z)=-B_{rs}(z)\varphi(s,B_{rs}(z)),
\end{equation}
where $\{\varphi(t,z)\}_{0<t<\infty}$ is a family of Herglotz functions such that 
$\varphi(t,z)$ is a measurable function of $t$ for each fixed $z$ 
(see Chapter 8 in Rosenblum--Rovnyak~\cite{RR} for further details). 
In what follows, we will write $B_t=B_{at}$ for short if no confusion occurs. 
Then, for any $t\geq 0$,
\[
k(t, z,w)=\frac{\overline{\varphi(t,B_t(w))}+\varphi(t,B_t(z))}{1-\overline{w} z}
\]
is a kernel function defined on $\D\times \D$ by Proposition \ref{prop:1-3}. 
Then, trivially, Assumptions (A) and (B) are valid. 
Moreover, Assumptions (C) and (D) follow from the following estimate:
\begin{equation*}
0\leq k(t,\lam,\lam)=\dfrac{2\Real\varphi(t,B_t(\lam))}{1-|\lam|^2}\leq 
\dfrac{2(1+|B_t(\lam)|)}{(1-|\lam|^2)(1-|B_t(\lam)|)}
\end{equation*}
for any $\lam$ in $\D$ (see (8--10) in Page 193 of Rosenblum--Rovnyak~\cite{RR}). 
Moreover, trivially, Assumption (E) is valid for $B(t,z)=B_t(z)$.  

\begin{lem}\label{lem:2-1}
Let $B_t(z)=B_{at}(z)$ $(t\geq a\geq 0)$ be the L\"{o}wner semigroup. Then, 
\[
\dfrac{\p}{\p t} \dfrac{1-\overline{B_{t}(\lam)}B_{t}(z)}{1-\overline{\lam}z}
=\dfrac{\overline{\varphi(t,B_{t}(\lam))}+\varphi(t,B_{t}(z))}{1-\overline{\lam}z}
\overline{B_{t}(\lam)}B_{t}(z).
\]
\end{lem}

\begin{proof}
By the nonlinear L\"{o}wner equation (\ref{eq:5-1}), 
we have the conclusion. 
\end{proof}

Let $\cH(B_t)$ denote the reproducing kernel Hilbert space 
generated by kernel function 
\[
\dfrac{1-\overline{B_t(w)}B_t(z)}{1-\overline{w}z},
\]
which is called the de Branges--Rovnyak space induced by $B_t$ (see de Branges--Rovnyak~\cite{dBR, dBR2}, 
Fricain--Mashreghi~\cite{FM} and Sarason~\cite{Sarason}). 
Every function in $\cH(B_t)$ has a Fourier type integral representation. 

\begin{thm}[Ghosechowdhury~\cite{G2}]\label{thm:5-2}
Let $B_t(z)=B_{at}(z)$ $(t\geq a\geq 0)$ be a L\"{o}wner semigroup. Then, 
\[
\cH(B_b)=\C \oplus \M(\bI_B; L^2(\cL(\varphi(t,B_t)))).
\]
\end{thm}

\begin{proof}
We will write $h_{\lam}(t,z)=\overline{B_t(\lam)}k_{\lam}(t,z)$ for any $\lam$ in $\D$. 
It follows from Theorem \ref{thm:2-2} and Lemma \ref{lem:2-1} that 
\begin{align*}
\la \bI_B h_{\lam}, \bI_B h_{\mu} \ra_{\M(\bI_B)}
&=\la  h_{\lam},  h_{\mu} \ra_{L^2(\cL(\varphi(t,B_{t})))}\\
&=\int_a^b \la \overline{B_t(\lam)}k_{\lam}(t,\cdot),\overline{B_t(\mu)}k_{\mu}(t,\cdot) \ra_{\cL(\varphi(t,B_{t}))}\ dt\\
&=\int_a^b \overline{B_t(\lam)}B_t(\mu)k(t,\mu,\lam)\ dt\\
&=\int_a^b \dfrac{\p}{\p t}\dfrac{1-\overline{B_t(\lam)}B_t(\mu)}{1-\overline{\lam}\mu}\ dt\\
&=\dfrac{1-\overline{B_b(\lam)}B_b(\mu)}{1-\overline{\lam}\mu}-1.
\end{align*}
Since $1$ belongs to $\cH(B_b)$, 
\[
\dfrac{1-\overline{B_b(\lam)}B_b(z)}{1-\overline{\lam}z}-1
\]
is the reproducing kernel of $\cH(B_b)\ominus \C$, where the orthogonal complement is taken with respect to the inner product of $\cH(B_b)$. 
Hence, we have that $\M(\bI_B)=\cH(B_b)\ominus \C$. 
This concludes the proof. 
\end{proof}

Here, we would like to note that 
Ghosechowdhury's theorem suggests a new insight to the problem 
\begin{center}
``\textit{Find concrete elements in de Branges--Rovnyak spaces}''
\end{center}
discussed in (II-8) of Sarason~\cite{Sarason}, Sections 18.6 and 23.8 of Fricain--Mashreghi~\cite{FM}, Fricain--Hartmann--Ross~\cite{FHR} and {\L}anucha--Nowak~\cite{LN}. 
Combining 
Theorem \ref{thm:5-2} with 
the observations in Example \ref{sec:1-1-6}, 
we see that every function in $\cH(B_{b})$ has a double integral representation. 
Indeed, let $F$ be any function in $\cH(B_b)$. 
Then, by Theorem \ref{thm:5-2}, 
there exists some function $f$ in $L^2(\cL(\varphi(t,B_{t})))$ such that  
\[
F(z)=F(0)+(\bI_B f)(z)=F(0)+\int_a^bB_t(z)f(t,z)\ dt. 
\]
Moreover, for almost all $t$, 
it follows from the observation in Example \ref{sec:1-1-6} 
that there exists some function $g$ in $L^2(\cL(\varphi_{\xi}(B_{t})))$ such that
\[
f(t,z)=\int_{\bT} g(\xi, t, z)\ d\mu_t(\xi).
\]
Hence, we have that 
\[
F(z)=F(0)+\int_a^bB_t(z)\left(\int_{\bT} g(\xi, t, z)\ d\mu_t(\xi)\right)\ dt. 
\]
Furthermore, 
since 
\[
\frac{\overline{\varphi_{\xi}(B_t(\lam))}+\varphi_{\xi}(B_t(z))}{1-\overline{\lam}z}
=\frac{1-\overline{B_t(\lam)}B_t(z)}{1-\overline{\lam}z}\cdot
\frac{2}{\left(1-\overline{\xi B_t(\lam)}\right)\left( 1-\xi B_t(z)\right)}
\]
by (\ref{eq:3-3}), 
\begin{equation}\label{eq:2-3}
F(z)=\int_a^b
B_t(z)\frac{1-\overline{B_t(\lam)}B_t(z)}{1-\overline{\lam}z}\left(\int_{\bT}\frac{h(\xi,t)}{\left(1-\overline{\xi B_t(\lam)}\right)\left(1-\xi B_t(z)\right)}\ d\mu_t(\xi)\right)\ dt
\end{equation}
belongs to $\cH(B_{b})$ for any bounded Borel function $h$. 
We shall give an answer to Question 6.6 in \cite{FHR}. 

\begin{cor}\label{cor:3-7}
Let $B_t(z)=B_{at}(z)$ $(t\geq a\geq 0)$ be the L\"{o}wner semigroup associated with the Koebe functions
\[
f(t,z)=\frac{e^tz}{(1-z)^2}\quad (t\geq 0).
\]
Then, 
\[
\log \frac{1-B_t(z)}{1-z}
\]
belongs to $\cH(B_t)$. 
\end{cor}

\begin{proof}
First, note that 
$B_t/z$ is outer by Theorem 3.17 in Duren~\cite{Duren}, 
and is extreme by Figure 8.2 in Rosenblum--Rovnyak~\cite{RR}. 
Let $b$ be an arbitrary point satisfying $a\leq b$. 
Since $\mu_t$ is the Dirac measure centered on $\xi=-1$ for all $t\geq 0$ 
and
\[
\varphi(t,z)=\varphi_{-1}(z)=\frac{1-z}{1+z}
\]
for the Koebe functions, 
it follows from (\ref{eq:2-3}) that  
\[
F(z)=\int_a^b 
B_t(z)\frac{1-\overline{B_t(\lam)}B_t(z)}{1-\overline{\lam}z}
\frac{h(t)}{\left(1+\overline{B_t(\lam)}\right)\left(1+B_t(z)\right)}\ dt
\]
belongs to $\cH(B_b)$ for any bounded Borel function $h$ and any $\lam$ in $\D$. 
Setting $h(t)=1$ and $\lam=0$, by the nonlinear L\"{o}wner equation (\ref{eq:5-1}), we have that 
\begin{align*}
F(z)
&=\int_a^b \frac{B_t(z)}{1+B_t(z)}\ dt\\
&=\int_a^b\frac{-1}{1-B_t(z)}\left(-B_t(z)\frac{1-B_t(z)}{1+B_t(z)}\right)\ dt\\
&=\int_a^b \frac{-1}{1-B_t(z)}\frac{\p B_t}{\p t}(z)\ dt\\
&=\int_a^b \frac{\p}{\p t}\log\left(1-B_t(z) \right)\ dt\\
&=\log \frac{1-B_b(z)}{1-z}.
\end{align*}
This concludes the proof. 
\end{proof}






\section{Pick spaces}\label{sec:6}

Let $\varphi$ be a Pick function, that is, 
$\varphi$ is a holomorphic self-map of the upper half-plane $\bH=\{z\in \C: \operatorname{Im}z>0\}$. 
Every Pick function has a Nevanlinna representation 
\begin{equation}\label{eq:5-0}
\varphi(z)=b+cz+\frac{1}{\pi}\int_{-\infty}^{\infty}\left( \frac{1}{t-z}-\frac{t}{1+t^2}\right)\ d\mu(t),
\end{equation}
where $b$ is real, $c\geq 0$ and $\mu$ is a nonnegative Borel measure on $(-\infty,\infty)$ satisfying
\[
\int_{-\infty}^{\infty}\frac{1}{1+t^2}\ d\mu(t)<\infty,
\]
and the converse is also true (see Theorem 5.3 in Rosenblum--Rovnyak~\cite{RR}). 

\begin{ex}\label{ex:5-1}\rm
Setting 
\[
\varphi_{\xi}(z)=\frac{1}{\xi -z}\quad (\xi \in \R),
\]
this $\varphi_{\xi}$ is a typical example of a Pick function, 
which is obtained from (\ref{eq:5-0}) if we choose $\mu$ as the Dirac measure centered on $\xi$. 
\end{ex}

Let $\varphi$ be a Pick function. 
Then, $\psi$ will denote the function defined by the following diagram:
\begin{equation}\label{eq:6-2}
\begin{CD}
\bH @>{\varphi}>> \bH \\
@A{T}AA    @VV{T^{-1}}V \\
\D  @>>{\psi}>  \D,
\end{CD}
\end{equation}
where 
\[
T(z)=i\dfrac{1+z}{1-z}\quad \text{and}\quad T^{-1}(w)=\frac{w-i}{w+i}.
\]
Then, $\psi$ is a holomorphic self-map of $\D$. 
Let $\cH(\psi)$ denote the de Branges--Rovnyak space induced by $\psi$, that is, 
$\cH(\psi)$ is the reproducing kernel Hilbert space generated by kernel function
\[
\frac{1-\overline{\psi(w)}\psi(z)}{1-\overline{w}z}.
\] 
Let $\alpha$ and $\beta$ be arbitrary two points in $\bH$. 
We set $\lam=T^{-1}(\alpha)$ and $\mu=T^{-1}(\beta)$. 
Then, by direct calculation, we have that 
\begin{equation}\label{eq:6-3}
\frac{\varphi(\beta)-\overline{\varphi(\alpha)}}{\beta-\overline{\alpha}}
=\frac{1-\overline{\lam}}{1-\overline{\psi(\lam)}} \cdot \frac{1-\mu}{1-\psi(\mu)} \cdot \frac{1-\overline{\psi(\lam)}\psi(\mu)}{1-\overline{\lam}\mu}.
\end{equation}
By (\ref{eq:6-3}) and Schur's theorem,
we have the following well-known fact. 
\begin{prop}\label{prop:5-2}
Let $\varphi$ be a Pick function. Then, 
\[
\frac{\varphi (z)-\overline{\varphi (w)}}{z-\overline{w}}
\]
is a kernel function on $\bH \times \bH$.
\end{prop}

Let $\cP(\varphi)$ denote
the reproducing kernel Hilbert space generated by the kernel function in Proposition \ref{prop:5-2}. 
We will call $\cP(\varphi)$ a Pick space.

\begin{prop}\label{thm:5-4}
Let $\varphi$ be a Pick function
with Nevanlinna representation (\ref{eq:5-0}), 
and let $\M$ denote the reproducing kernel Hilbert space over $\bH$ generated by kernel function 
\[
k(z,w)=\frac{1}{\pi}\int_{-\infty}^{\infty}\frac{1}{(t-\overline{w})(t-z)}\ d\mu(t).
\]
Then, 
\[
\cP(\varphi)=
\begin{cases}
\C + \M & (c \neq 0)\\
\M & (c=0 )
\end{cases}
\] 
as reproducing kernel Hilbert spaces,  
where we consider $\C$ to be the reproducing kernel Hilbert space consisting of constant functions on $\bH$. 
\end{prop}

\begin{proof}
It follows from the Nevanlinna representation (\ref{eq:5-0}) that 
\begin{align*}
\varphi(z)-\overline{\varphi(w)}
&=c(z-\overline{w})+\frac{1}{\pi}\int_{-\infty}^{\infty}\left( \frac{1}{t-z}-\frac{1}{t-\overline{w}}\right)\ d\mu(t)\\
&=c(z-\overline{w})+\frac{1}{\pi}\int_{-\infty}^{\infty}\frac{z-\overline{w}}{(t-\overline{w})(t-z)}\ d\mu(t).
\end{align*}
Hence, we have that 
\begin{equation*}
\frac{\varphi(z)-\overline{\varphi(w)}}{z-\overline{w}}
=c+\frac{1}{\pi}\int_{-\infty}^{\infty}\frac{1}{(t-\overline{w})(t-z)}\ d\mu(t).
\end{equation*}
By this identity and Aronszajn's sums of kernels theorem  
(see Theorem 5.4 in Paulsen--Raghupathi~\cite{PR}), we conclude the proof.
\end{proof}

Next, we shall give an isomorphism between $\cP(\varphi)$ and $\cH(\psi)$. 
\begin{thm}\label{lem:6-3}
Let $\varphi$ be a Pick function and $\psi$ be the function defined by (\ref{eq:6-2}). 
Then, 
\[
V: f(z)\mapsto \frac{1-T^{-1}(w)}{1-\psi\circ T^{-1}(w)}f\circ T^{-1}(w)\quad (w=T(z))
\]
gives an isometry from $\cH(\psi)$ onto $\cP(\varphi)$.   
\end{thm}

\begin{proof}
In this proof, we will write
\[
\frac{1-T^{-1}(w)}{1-\psi\circ T^{-1}(w)}f\circ T^{-1}(w)=
\frac{1-z}{1-\psi(z)}f(z)
\]
for short. 
It follows from (\ref{eq:6-3}) that 
\[
\frac{1-z}{1-\psi(z)}\cdot \dfrac{1-\overline{\psi(\lam)}\psi(z)}{1-\overline{\lam}z}
=\frac{1-\overline{\psi(\lam)}}{1-\overline{\lam}}
\cdot\frac{\varphi(w)-\overline{\varphi(\alpha)}}{w-\overline{\alpha}}
\]
belongs to $\cP(\varphi)$. 
Moreover, by (\ref{eq:6-3}) again, we have that
\begin{align*}
&\left\la V\dfrac{1-\overline{\psi(\lam)}\psi(z)}{1-\overline{\lam}z}, 
V\dfrac{1-\overline{\psi(\mu)}\psi(z)}{1-\overline{\mu}z} \right\ra_{\cP(\varphi)}\\
&\quad=\left\la \frac{1-z}{1-\psi(z)}\cdot \dfrac{1-\overline{\psi(\lam)}\psi(z)}{1-\overline{\lam}z}, 
\frac{1-z}{1-\psi(z)}\cdot\dfrac{1-\overline{\psi(\mu)}\psi(z)}{1-\overline{\mu}z} \right\ra_{\cP(\varphi)}\\
&\quad=\left\la\frac{1-\overline{\psi(\lam)}}{1-\overline{\lam}}
\cdot\frac{\varphi(w)-\overline{\varphi(\alpha)}}{w-\overline{\alpha}},
\frac{1-\overline{\psi(\mu)}}{1-\overline{\mu}}\cdot \frac{\varphi(w)-\overline{\varphi(\beta)}}{w-\overline{\beta}}\right\ra_{\cP(\varphi)}\\
&\quad= \frac{1-\overline{\psi(\lam)}}{1-\overline{\lam}}\cdot \frac{1-\psi(\mu)}{1-\mu}
\cdot\frac{\varphi(\beta)-\overline{\varphi(\alpha)}}{\beta-\overline{\alpha}}\\
&\quad=\frac{1-\overline{\psi(\lam)}\psi(\mu)}{1-\overline{\lam}\mu}\\
&\quad=\left\la \dfrac{1-\overline{\psi(\lam)}\psi(z)}{1-\overline{\lam}z}, 
\dfrac{1-\overline{\psi(\mu)}\psi(z)}{1-\overline{\mu}z} \right\ra_{\cH(\psi)}.
\end{align*}
Thus, we have the conclusion.
\end{proof}

As an application of Theorem \ref{lem:6-3}, 
nontrivial elements in $\cH(\psi)$ are obtained. 
For example, if $c\neq 0$ in Proposition \ref{thm:5-4}, 
then constant function $1$ belongs to $\cP(\varphi)$. 
Hence, by Theorem \ref{lem:6-3}, we have the following corollary.
\begin{cor}\label{cor:5-5}
If $c\neq 0$, then
\[
\frac{1-\psi(z)}{1-z}
\]
belongs to $\cH(\psi)$.
\end{cor}
In the next section, we will see that 
further elements in $\cH(\psi)$ will be obtained from analysis of $\cP(\varphi)$. 

\section{A chordal L\"{o}wner expansion}\label{sec:2-1-6}


We shall begin with the following fundamental example in the chordal L\"{o}wner theory. 
\begin{ex}\label{ex:6-1} \rm
We set 
\[
B_{rs}(z)=\sqrt{z^2-2(s-r)}\quad (0\leq r \leq s).
\] 
Then, $B_{rs}$ is a univalent (i.e. holomorphic and injective) map from $\bH$ onto $\bH \setminus \{iy: 0<y \leq \sqrt{2(r-s)}\}$. 
Moreover,  
$B_{rs}$ is the solution of  
\begin{equation}\label{eq:6-1}
\frac{\p B_{rs}}{\p s}(z)=-\frac{1}{B_{rs}(z)}
\end{equation}
with initial condition $B_{ss}(z)=z$. 
The differential equation (\ref{eq:6-1}) is the basic case of the chordal L\"{o}wner equation
\begin{equation}\label{eq:6-2-2}
\frac{\p B_{rs}}{\p s}(z)=\int_{\R}\frac{1}{\xi-B_{rs}(z)}\ d\nu_s(\xi)=\int_{\R}\varphi_{\xi}(B_{rs}(z))\ d\nu_s(\xi),
\end{equation}
where $\nu_s$ is a Borel measure on $\R$ and $\varphi_{\xi}$ is the function in Example \ref{ex:5-1}.
Indeed, choosing $\nu_s$ as the Dirac measure centered on $\xi=0$ for all $s\geq 0$ in (\ref{eq:6-2-2}), 
we have (\ref{eq:6-1}). 
\end{ex}
In this paper, we need only the above example.  
For further details of the general chordal L\"{o}wner equation, see Schlei{\ss}inger~\cite{Sch}. 
Let $B_t(z)=B_{at}(z)$ $(t\geq a\geq 0)$ be the family of functions in Example \ref{ex:6-1}.
Then, applying Theorem \ref{thm:2-2}, we consider the integral operator defined as follows: 
\[
\bI_{1/B}: 
f(t) \mapsto F(z)=\int_a^b \frac{1}{B_t(z)}f(t)\ dt \quad (f\in L^2([a,b])).
\]

\begin{lem}\label{lem:6-2}
Let $B_t(z)=B_{at}(z)$ $(t\geq a\geq 0)$ be the family of functions in Example \ref{ex:6-1}.
Then, for any $\alpha$ in $\bH$, 
\[
\frac{\p}{\p t}
\frac{B_t(z)-\overline{B_t(\alpha)}}{z-\overline{\alpha}}
=\frac{1}{\overline{B_t(\alpha)}B_t(z)}\cdot \frac{B_t(z)-\overline{B_t(\alpha)}}{z-\overline{\alpha}}.
\]
\end{lem}

\begin{proof}
By the chordal L\"{o}wner equation (\ref{eq:6-1}), we have the conclusion. 
\end{proof}

\begin{thm}\label{thm:6-3}
Let $B_t(z)=B_{at}(z)$ $(t\geq a\geq 0)$ be the family of functions in Example \ref{ex:6-1}.
Then, 
\[
\cP(B_b)=\exp \M(\bI_{1/B}; L^2([a,b])).
\]
\end{thm}

\begin{proof}
We set $h_{\alpha}(t)=1/\overline{B_t(\alpha)}$ for any $\alpha$ in $\bH$. 
Then, $h_{\alpha}$ belongs to $L^2([a,b])$ and
\[
\exp ((\bI_{1/B} h_{\alpha})(z))=\exp\left(\int_a^b \frac{1}{\overline{B_t(\alpha)}B_t(z)}\ dt\right)
\]
is the kernel function generating $\exp \M(\bI_{1/B})$ 
(see Exercise (k) in p. 320 of Nikolski~\cite{Nik} or Chapter 7 in Paulsen--Raghupathi~\cite{PR}). 
Moreover, by Lemma \ref{lem:6-2}, we have that  
\begin{align*}
\exp ((\bI_{1/B} h_{\alpha})(z))
&=\exp \left(\int_a^b \frac{1}{\overline{B_t(\alpha)}B_t(z)}\ dt\right)\\
&=\exp \left(\int_a^b\frac{w-\overline{\alpha}}{B_t(w)-\overline{B_t(\alpha)}}
\cdot \dfrac{\p}{\p t}\dfrac{B_t(w)-\overline{B_t(\alpha)}}{w-\overline{\alpha}}\ dt\right)
\\
&=\exp \left(\int_a^b \dfrac{\p}{\p t}\log \left(\dfrac{B_t(w)-\overline{B_t(\alpha)}}{w-\overline{\alpha}}\right)\ dt\right)\\
&=\frac{B_b(w)-\overline{B_b(\alpha)}}{w-\overline{\alpha}}.
\end{align*} 
This concludes the proof.
\end{proof}




\begin{cor}\label{cor:6-4}
Let $B_t(z)=B_{at}(z)$ $(t\geq a\geq 0)$ be the family of functions in Example \ref{ex:6-1}.
Then, 
\[
\exp(z-B_t(z))
\]
belongs to $\cP(B_t)$. 
\end{cor}

\begin{proof}
It follows from Theorem \ref{thm:6-3} that  
\[
F(z)=\exp \left(\int_a^b \frac{1}{B_t(z)}f(t)\ dt\right)
\]
belongs to $\cP(B_b)$ for any function $f$ in $L^2([a,b])$. 
Setting $f(t)=1$, by the chordal L\"{o}wner equation (\ref{eq:6-1}), we have that 
\begin{align*}
F(z)
&=\exp \left(\int_a^b \frac{1}{B_t(z)}\ dt\right)\\
&=\exp \left(-\int_a^b \frac{\p B_t}{\p t}(z)\ dt\right)\\
&=\exp(z-B_b(z))
\end{align*}
This concludes the proof. 
\end{proof}


\begin{rem}\rm 
In the same way as the proof of Theorem \ref{thm:6-3}, 
general chordal L\"{o}wner expansions analogous to Theorem \ref{thm:5-2}  are obtained from 
the chordal L\"{o}wner equation (\ref{eq:6-2-2}). 
\end{rem}


\begin{acknowledgment}\rm
The author would like to thank 
Professor Ikkei Hotta (Yamaguchi University) and Dr. Takuya Murayama (Kobe University). 
This research is greatly indebted to their unpublished surveys on the L\"{o}wner theory. 
This work was supported by JSPS KAKENHI Grant Number JP24K06771 and the Research Institute for Mathematical Sciences,
an International Joint Usage/Research Center located in Kyoto University.
\end{acknowledgment}


\begin{thebibliography}{References}

\bibitem{Ando}T. Ando, 
\textit{De Branges spaces and analytic operator functions}. Lecture Notes, Hokkaido University, Sapporo, 1990. 

\bibitem{Conway}
J. B. Conway, 
\textit{Functions of one complex variable}. II.
Grad. Texts in Math., 159
Springer-Verlag, New York, 1995. 

\bibitem{de Branges2}
L. de Branges, 
\textit{L\"{o}wner expansions}.
J. Math. Anal. Appl. 100 (1984), no. 1, 323--337.

\bibitem{dBR}
L. de Branges and J. Rovnyak, 
\textit{Appendix on square summable power series} in  
\textit{Canonical models in quantum scattering theory}. 
Perturbation Theory and its Applications in Quantum Mechanics (Proc. Adv. Sem. Math. Res. Center, U.S. Army, Theoret. Chem. Inst., Univ. of Wisconsin, Madison, Wis., 1965), pp. 295--392
John Wiley \& Sons, Inc., New York--London--Sydney, 1966

\bibitem{dBR2}
L. de Branges and J. Rovnyak, 
\textit{Square summable power series}.
Holt, Rinehart and Winston, New York--Toronto--London, 1966.

\bibitem{Duren}
P. L. Duren, 
\textit{Theory of $H^p$ spaces}.
Pure Appl. Math., Vol. 38
Academic Press, New York--London, 1970. 

\bibitem{FHR}
E. Fricain, A. Hartmann and W. T. Ross, 
\textit{Concrete examples of $\cH(b)$ spaces}.
Comput. Methods Funct. Theory 16 (2016), no. 2, 287--306.

\bibitem{FM}
E. Fricain and J. Mashreghi, 
\textit{The theory of $\cH(b)$ spaces}. Vol. 2.
New Math. Monogr., 21
Cambridge University Press, Cambridge, 2016. 

\bibitem{G1}
S. Ghosechowdhury, 
\textit{An expansion theorem for state space of unitary linear system whose transfer function is a Riemann mapping function}.
Reproducing kernels and their applications (Newark, DE, 1997), 81--95.
Int. Soc. Anal. Appl. Comput., 3 
Kluwer Academic Publishers, Dordrecht, 1999

\bibitem{G2}
S. Ghosechowdhury, 
\textit{L\"{o}wner expansions}. 
Math. Nachr. 210 (2000), pp. 111--126.

\bibitem{LN}
B. {\L}anucha and M. T. Nowak,
\textit{Examples of de Branges--Rovnyak spaces generated by nonextreme functions}.
Ann. Acad. Sci. Fenn. Math. 44 (2019), no. 1, 449--457.

\bibitem{Nik}
N. K. Nikolski,
\textit{Operators, functions, and systems: an easy reading}. Vol. 1. 
Hardy, Hankel, and Toeplitz. Translated from the French by Andreas Hartmann. Mathematical Surveys and Monographs, 
\textbf{92}. American Mathematical Society, Providence, RI, 2002. 

\bibitem{PR}
V. I. Paulsen and M. Raghupathi, 
\textit{An introduction to the theory of reproducing kernel Hilbert spaces}.
Cambridge Studies in Advanced Mathematics, 152. Cambridge University Press, Cambridge, 2016. 

\bibitem{RR}M. Rosenblum and J. Rovnyak, 
\textit{Topics in Hardy Classes and Univalent Functions}. 
Birkh\"{a}user Verlag, 1994.

\bibitem{Sarason}D. Sarason, 
\textit{Sub-Hardy Hilbert spaces in the unit disk}. 
University of Arkansas Lecture Notes in the Mathematical Sciences, 10. 
A Wiley-Interscience Publication. John Wiley \& Sons, Inc., New York, 1994. 

\bibitem{Sch}
S. Schlei{\ss}inger, 
\textit{The chordal Loewner equation and monotone probability theory}. 
Infin. Dimens. Anal. Quantum Probab. Relat. Top. 20 (2017), no. 3, 1750016.  

\bibitem{S}
L.\ Shulman, 
\textit{Perturbations of unitary transformations}.
Amer. J. Math.   91 (1969), 267--288.

\bibitem{VN1}V. I. Vasyunin and N. K. Nikol'ski\u{\i}, 
\textit{Quasiorthogonal decompositions with respect to complementary metrics, and estimates of univalent functions}. 
Leningrad Math.\ J. 2 (1991), no. 4, pp. 691--764. 

\bibitem{VN2}
V. I. Vasyunin and N. K. Nikol'ski\u{\i}, 
\textit{Operator-valued measures and coefficients of univalent functions}.
St. Petersburg Math. J. 3 (1992), no. 6, pp. 1199--1270.


\end{thebibliography}


\end{document}


